\documentclass{exos}
\usepackage{main}

\title{Activité : Variations de fonctions}
\date{Chapitre 8 Partie 1}
\author{Seconde}

\begin{document}
\maketitle

Au cours d'une expérience, la température d'une pièce change pour stimuler une réaction chimique. L'expérience dure \qty{24}{\hour}, et on note $T(t)$ la température (en \unit{\celsius}) de la pièce en fonction du temps $t$ (en \unit{\hour}). La courbe représentative de la fonction $T$ est représentée sur le repère ci-après :

\begin{center}
\begin{tikzpicture}
\begin{axis}[xmin=0,xmax=24,ymin=-30,ymax=55,
  enlargelimits, 
  xtick={0,1,...,24}, 
  width=12.5cm,
  xlabel=$t$ (en \unit{\hour}),
  ylabel=$T$ (en \unit{\celsius})]
\addplot[smooth] coordinates {
(0,25)(4,-20)(16,50)(24,25)
};
\draw (0,25) node {$\bullet$};
\draw (24,25) node {$\bullet$};
\end{axis}
\end{tikzpicture}
\end{center}
\begin{alphaquestions}
\item Quelle est la température de la salle au bout de \qty{4}{\hour} ?
\item L'expérience se déroule en trois étapes :
\begin{enumerate}
  \item Ajout de l'azote liquide froid pendant un temps long.
  \item Réaction exothermique.
  \item Phase de repos.
\end{enumerate}
À l'aide de la représentation de $T$, donner l'heure de départ et de fin de chacune de ces étapes.
\item Ordonner les heures suivantes : $10; 7; 15; 5$. Ordonner les températures correspondantes $T(10); T(7); T(15)$ et $T(5)$. Que remarquez-vous ?
\item Même question que la question précédente, mais pour les heures $23; 18; 21$.
\item On dit que $T$ est croissante sur $[4;17]$ et décroissante sur $[17;24]$. Que peut-on dire de $T$ sur l'intervalle $[0;4]$ ?
\end{alphaquestions}
\newpage
Le lendemain, on tente la même expérience dans une pièce différente. On obtient le résultat représenté par la courbe suivante :

\begin{center}
\begin{tikzpicture}
\begin{axis}[xmin=0,xmax=24,ymin=-30,ymax=55,
  enlargelimits, 
  xtick={0,1,...,24}, 
  width=12.5cm,
  xlabel=$t$ (en \unit{\hour}),
  ylabel=$T$ (en \unit{\celsius})]
\addplot[smooth] coordinates {
(0,18)(8,-12)(17,25)(19,-5)(24,18)
};
\draw (0,18) node {$\bullet$};
\draw (24,18) node {$\bullet$};
\end{axis}
\end{tikzpicture}
\end{center}
S'inspirer de la phrase \og La fonction $T$ est \textbf{croissante} sur $[19;24]$ \fg pour compléter les suivantes :
\begin{alphaquestions}
\item La fonction $T$ est \dots sur $[0;8]$.
\item La fonction $T$ est décroissante sur $\dots$.
\item La fonction $T$ est croissante sur $\dots$.
\end{alphaquestions}

\end{document}